\documentclass[10pt]{article}
% Shared preamble of the problems/MHD/setup_notes documents.
% Each case document does % Shared preamble of the problems/MHD/setup_notes documents.
% Each case document does % Shared preamble of the problems/MHD/setup_notes documents.
% Each case document does \input{preamble} right after \documentclass.
\usepackage[margin=0.85in]{geometry}
\usepackage[T1]{fontenc}
\usepackage{lmodern}
\usepackage{amsmath,amssymb,bm}
\usepackage{array,booktabs,tabularx,longtable,multirow}
\usepackage{xcolor}
\usepackage{enumitem}
\usepackage[normalem]{ulem}
\usepackage[hyphens]{url}
\usepackage[colorlinks=true,linkcolor=blue!50!black,urlcolor=blue!50!black]{hyperref}
\usepackage{caption}
\usepackage{graphicx}
\setlist{itemsep=1pt,topsep=3pt}
\sloppy
\emergencystretch=3em
\newcolumntype{Y}{>{\raggedright\arraybackslash}X}
\newcolumntype{L}[1]{>{\raggedright\arraybackslash}p{#1}}

% \code{...}: typewriter with breakable underscores (write \_ in the argument).
\makeatletter
\newcommand{\code}[1]{{\def\_{\char`\_\discretionary{}{}{}}\texttt{#1}}}
\makeatother

% --- Common-item marks -------------------------------------------------------
% \cm{x}  solid underline : identical in EVERY case to which the item applies
% \mj{x}  dashed underline: shared by a MAJORITY of those cases, but not all
\newcommand{\cm}[1]{\uline{#1}}
\newcommand{\mj}[1]{\dashuline{#1}}

\newcommand{\commonlegend}{%
\begin{center}
\fcolorbox{black!40}{black!3}{\parbox{0.92\textwidth}{\small
\textbf{How to read the underlines.}
\cm{Solid underline}: the setting (or the code path) is \emph{identical in every case
of \code{problems/MHD} to which it applies} --- all nine cases for general settings, the eight
\code{DSGS\_MHD} cases for DynSGS settings, the three \code{:positivity\_method => "conservative"}
cases for the element limiter.
\mj{Dashed underline}: shared by a \emph{majority} of those cases but not all of them.
No underline: specific to this case or to a minority.
The cross-case comparison is in \code{summary.pdf} of this directory.}}
\end{center}}

% --- Notation ----------------------------------------------------------------
\newcommand{\bq}{\mathbf{q}}
\newcommand{\bv}{\mathbf{v}}
\newcommand{\bmom}{\mathbf{m}}
\newcommand{\bB}{\mathbf{B}}
\newcommand{\bF}{\mathbf{F}}
\newcommand{\Et}{\tilde E}
\newcommand{\eps}{\varepsilon}
\newcommand{\dsgs}{DynSGS}
\newcommand{\CR}{C_R}
\newcommand{\Cmax}{C_{\max}}
\newcommand{\Cmin}{C_{\min}}
\newcommand{\nuc}{\nu_{c}}
\newcommand{\nufl}{\nu_{\mathrm{fl}}}
\newcommand{\commit}{\code{dcc3653}}
\newcommand{\branch}{\code{sm/MHD}}

% Two-column "setting" table: key | value as used | where / remark
\newenvironment{keytable}{%
\par\smallskip\noindent\small
\begin{longtable}{@{}L{0.24\textwidth}L{0.38\textwidth}L{0.32\textwidth}@{}}
\toprule
\textbf{Deck key / item} & \textbf{Value as run} & \textbf{Source / remark}\\
\midrule\endhead
\bottomrule\endfoot}{%
\end{longtable}\par}
 right after \documentclass.
\usepackage[margin=0.85in]{geometry}
\usepackage[T1]{fontenc}
\usepackage{lmodern}
\usepackage{amsmath,amssymb,bm}
\usepackage{array,booktabs,tabularx,longtable,multirow}
\usepackage{xcolor}
\usepackage{enumitem}
\usepackage[normalem]{ulem}
\usepackage[hyphens]{url}
\usepackage[colorlinks=true,linkcolor=blue!50!black,urlcolor=blue!50!black]{hyperref}
\usepackage{caption}
\usepackage{graphicx}
\setlist{itemsep=1pt,topsep=3pt}
\sloppy
\emergencystretch=3em
\newcolumntype{Y}{>{\raggedright\arraybackslash}X}
\newcolumntype{L}[1]{>{\raggedright\arraybackslash}p{#1}}

% \code{...}: typewriter with breakable underscores (write \_ in the argument).
\makeatletter
\newcommand{\code}[1]{{\def\_{\char`\_\discretionary{}{}{}}\texttt{#1}}}
\makeatother

% --- Common-item marks -------------------------------------------------------
% \cm{x}  solid underline : identical in EVERY case to which the item applies
% \mj{x}  dashed underline: shared by a MAJORITY of those cases, but not all
\newcommand{\cm}[1]{\uline{#1}}
\newcommand{\mj}[1]{\dashuline{#1}}

\newcommand{\commonlegend}{%
\begin{center}
\fcolorbox{black!40}{black!3}{\parbox{0.92\textwidth}{\small
\textbf{How to read the underlines.}
\cm{Solid underline}: the setting (or the code path) is \emph{identical in every case
of \code{problems/MHD} to which it applies} --- all nine cases for general settings, the eight
\code{DSGS\_MHD} cases for DynSGS settings, the three \code{:positivity\_method => "conservative"}
cases for the element limiter.
\mj{Dashed underline}: shared by a \emph{majority} of those cases but not all of them.
No underline: specific to this case or to a minority.
The cross-case comparison is in \code{summary.pdf} of this directory.}}
\end{center}}

% --- Notation ----------------------------------------------------------------
\newcommand{\bq}{\mathbf{q}}
\newcommand{\bv}{\mathbf{v}}
\newcommand{\bmom}{\mathbf{m}}
\newcommand{\bB}{\mathbf{B}}
\newcommand{\bF}{\mathbf{F}}
\newcommand{\Et}{\tilde E}
\newcommand{\eps}{\varepsilon}
\newcommand{\dsgs}{DynSGS}
\newcommand{\CR}{C_R}
\newcommand{\Cmax}{C_{\max}}
\newcommand{\Cmin}{C_{\min}}
\newcommand{\nuc}{\nu_{c}}
\newcommand{\nufl}{\nu_{\mathrm{fl}}}
\newcommand{\commit}{\code{dcc3653}}
\newcommand{\branch}{\code{sm/MHD}}

% Two-column "setting" table: key | value as used | where / remark
\newenvironment{keytable}{%
\par\smallskip\noindent\small
\begin{longtable}{@{}L{0.24\textwidth}L{0.38\textwidth}L{0.32\textwidth}@{}}
\toprule
\textbf{Deck key / item} & \textbf{Value as run} & \textbf{Source / remark}\\
\midrule\endhead
\bottomrule\endfoot}{%
\end{longtable}\par}
 right after \documentclass.
\usepackage[margin=0.85in]{geometry}
\usepackage[T1]{fontenc}
\usepackage{lmodern}
\usepackage{amsmath,amssymb,bm}
\usepackage{array,booktabs,tabularx,longtable,multirow}
\usepackage{xcolor}
\usepackage{enumitem}
\usepackage[normalem]{ulem}
\usepackage[hyphens]{url}
\usepackage[colorlinks=true,linkcolor=blue!50!black,urlcolor=blue!50!black]{hyperref}
\usepackage{caption}
\usepackage{graphicx}
\setlist{itemsep=1pt,topsep=3pt}
\sloppy
\emergencystretch=3em
\newcolumntype{Y}{>{\raggedright\arraybackslash}X}
\newcolumntype{L}[1]{>{\raggedright\arraybackslash}p{#1}}

% \code{...}: typewriter with breakable underscores (write \_ in the argument).
\makeatletter
\newcommand{\code}[1]{{\def\_{\char`\_\discretionary{}{}{}}\texttt{#1}}}
\makeatother

% --- Common-item marks -------------------------------------------------------
% \cm{x}  solid underline : identical in EVERY case to which the item applies
% \mj{x}  dashed underline: shared by a MAJORITY of those cases, but not all
\newcommand{\cm}[1]{\uline{#1}}
\newcommand{\mj}[1]{\dashuline{#1}}

\newcommand{\commonlegend}{%
\begin{center}
\fcolorbox{black!40}{black!3}{\parbox{0.92\textwidth}{\small
\textbf{How to read the underlines.}
\cm{Solid underline}: the setting (or the code path) is \emph{identical in every case
of \code{problems/MHD} to which it applies} --- all nine cases for general settings, the eight
\code{DSGS\_MHD} cases for DynSGS settings, the three \code{:positivity\_method => "conservative"}
cases for the element limiter.
\mj{Dashed underline}: shared by a \emph{majority} of those cases but not all of them.
No underline: specific to this case or to a minority.
The cross-case comparison is in \code{summary.pdf} of this directory.}}
\end{center}}

% --- Notation ----------------------------------------------------------------
\newcommand{\bq}{\mathbf{q}}
\newcommand{\bv}{\mathbf{v}}
\newcommand{\bmom}{\mathbf{m}}
\newcommand{\bB}{\mathbf{B}}
\newcommand{\bF}{\mathbf{F}}
\newcommand{\Et}{\tilde E}
\newcommand{\eps}{\varepsilon}
\newcommand{\dsgs}{DynSGS}
\newcommand{\CR}{C_R}
\newcommand{\Cmax}{C_{\max}}
\newcommand{\Cmin}{C_{\min}}
\newcommand{\nuc}{\nu_{c}}
\newcommand{\nufl}{\nu_{\mathrm{fl}}}
\newcommand{\commit}{\code{dcc3653}}
\newcommand{\branch}{\code{sm/MHD}}

% Two-column "setting" table: key | value as used | where / remark
\newenvironment{keytable}{%
\par\smallskip\noindent\small
\begin{longtable}{@{}L{0.24\textwidth}L{0.38\textwidth}L{0.32\textwidth}@{}}
\toprule
\textbf{Deck key / item} & \textbf{Value as run} & \textbf{Source / remark}\\
\midrule\endhead
\bottomrule\endfoot}{%
\end{longtable}\par}

\title{\code{problems/MHD/fluxEmergenceSon2025}\\[2pt]
\large DynSGS and positivity setup, as coded on branch \branch{} (commit \commit)}
\author{}\date{}
\begin{document}
\maketitle
\commonlegend

\noindent\textbf{Scope.} Read off \code{user\_inputs.jl}, the case files and the kernels they reach on
\branch{}. This case pairs the conserved-form \dsgs{} on the \emph{departure from a magnetostatic
reference state} with a \emph{case-local, non-conservative} density/pressure floor applied as the
Runge--Kutta stage limiter. Its sibling \code{fluxEmergenceSon2025DSGS} drops the floor.

\section{The problem as set up}
\begin{keytable}
Equations & 2D ideal \cm{GLM-MHD} with gravity $g=1/\gamma$, \cm{state $(\rho,\rho u,\rho v,E,\rho w,B_x,B_y,B_z,\psi)$} & Son, Jang \& Magara (2025)\\
EOS & \cm{$p=(\gamma-1)(E-\tfrac12|\bmom|^2/\rho-\tfrac12|\bB|^2-\tfrac12\psi^2)$}, $\gamma=1.05$ & \code{user\_flux.jl}\\
Domain, mesh & $[0,80]\times[0,35]\,H_0$, periodic in $x$, walls \code{bottom}/\code{top}; \code{FE\_80x35.msh}, $80\times35$ quads of $1\,H_0$ & no refinement\\
Discretization & \cm{CG spectral elements, LGL nodes, lumped mass}; \mj{$N=4$} & \\
Time integration & \cm{\code{CarpenterKennedy2N54}} \emph{with} \code{stage\_limiter! = fe\_positivity\_limiter!}; \cm{fixed $\Delta t$} $=7.5\times10^{-3}\tau_0$, $t_{end}=54\,\tau_0$ & \code{user\_inputs.jl:94, 101}\\
Variables & \cm{\code{TOTAL()}, \code{"energy"}} & \\
Reference state $q_e(z)$ & two-temperature magnetostatic atmosphere ($T_{ch}=1$, $T_{cor}=25$, $z_{cor}=18$, $w_{tr}=0.6$) with a flux sheet $B_x(z)$ on $[4,8]\,H_0$, $\beta_*=1$; $\bv_e=0$ & \code{initialize.jl}; \code{qe[:,end]} $=p_e$\\
Initial data & $q_e$ plus $u=f(z)\,0.05\sin(2\pi(x-40)/20)$ for $|x-40|<10$ & paper eq.~5\\
Sources & \mj{\code{:lsource => true}}: gravity ($-\rho g$, $-\rho v g$), well-balanced correction of the discrete reference balance (\code{fe\_well\_balanced = true}), $\psi$ damping, absorbing layer $z>30$: $-\sigma(z)(q-q_e)$, $\sigma=2\sin^2\!\big(\tfrac\pi2\tfrac{z-30}{z_{top}-30}\big)$ & \code{user\_source.jl}\\
GLM cleaning & \cm{$c_h=\max(|\bv|+c_f)$ of the initial data, constant}; $S_\psi=-(\alpha_pc_h/\Delta h)\psi$, $\alpha_p=0.2$, $\Delta h$ the smallest nodal spacing & paper eq.~12\\
Boundary conditions & walls: Dirichlet normal momentum $=0$; \code{bottom} also normal $B=0$ & \code{user\_bc.jl}\\
Other switches & \cm{\code{:lfilter}, \code{:lkep}, \code{:entropy\_variables}, \code{:lrichardson}, \code{:ladapt} false} & \\
\end{keytable}

\section{DynSGS: deck keys and effective values}
\begin{keytable}
\code{:visc\_model} & \cm{\code{DSGS\_MHD()}}, \code{:lvisc => true} & \code{SGS.jl:1598}\\
\code{:}$\mu$ & \cm{all 1} & \\
\code{:dsgs\_sensor} & \mj{\code{"legacy"}} & \\
\code{:dsgs\_CR}, \code{:dsgs\_Cmax} & \mj{1.0}, \mj{0.5} & \\
\code{:dsgs\_Cmin} & 0.03 & background floor $0.03\,\Delta w$\\
\code{:dsgs\_gamma} & \cm{$=\gamma_{mhd}$} $=1.05$ & \\
\code{:dsgs\_Prt} & \cm{1.0} & enters $\gamma(\gamma-1)/\Pr=0.0525$\\
\code{:dsgs\_nazarov\_energy} & true & split energy flux (below)\\
\code{:dsgs\_norms} & \mj{\code{"element"}} & \\
\code{:dsgs\_local\_rel} & \cm{1.0} & \\
\code{:dsgs\_conserved} & \mj{true} & \\
\code{:dsgs\_nodal\_rho} & true in the deck, \cm{inactive} & \code{rhs.jl:1526} forces it off under \code{:dsgs\_conserved}\\
\code{:dsgs\_ref\_weight} & \mj{false} (default) & \\
\code{:dsgs\_fast\_floors}, \code{:dsgs\_thermal\_residual} & \mj{false} (defaults) & \\
\code{:dsgs\_rel} & \cm{1.0} (default) & not used by element norms\\
\code{:dsgs\_hold\_steps} & \mj{2} (default) & \\
\code{:dsgs\_cutoff}, \code{:dsgs\_freeze\_stage}, \code{:dsgs\_reference} & \cm{0}, \cm{false}, \cm{false} (defaults) & \\
\code{:ldsgs\_nodal} & \cm{false} (element form) & \\
\end{keytable}

\section{DynSGS as it runs for this case}
$\Delta^{elem}_K=1\,H_0$, \cm{$\Delta_K=\Delta^{elem}_K/(N+1)$} $=0.2\,H_0$, $h=\Delta t$.
\begin{enumerate}[leftmargin=*]
\item \textbf{History and hold.} \cm{Rotation once per step}; \mj{two held steps}: $R\equiv0$, so
$\nu=\nufl=0.03\,\Delta_Kw_K$ during steps 1--2.
\item \textbf{Legacy residual}, every stage:
$R^s_I=\big|\tfrac{3q^s_I-4q^{n,s}_I+q^{n-1,s}_I}{2h}-(M^{-1}\mathcal R_{inv})^s_I\big|$, with the inviscid RHS
(fluxes and \emph{all} sources: gravity, well-balanced term, $\psi$ damping, absorbing layer) assembled over
ranks and the periodic seam. The Dirichlet-node zeroing exists only on the residual-sensor path, so the
wall nodes are \emph{not} zeroed here.
\item \textbf{Element normalization} (\code{SGS.jl:1716--1761}): over the $25$ nodes of $K$,
$\langle q\rangle_K$ (plain average), $S^s_K=\max|q^s-\langle q^s\rangle_K|$,
$\bar\rho_K$, $\bar p_K$ (without $\psi$), $\bar c_K=\sqrt{\gamma\bar p_K/\bar\rho_K}$ (sound speed),
\[
n^s_K=\max\big(S^s_K,\ r_{loc}\,\sigma^s_K\big)+10^{-16},\quad
\sigma^\rho_K=\bar\rho_K,\ \sigma^{\bmom}_K=\bar\rho_K\bar c_K,\ \sigma^E_K=\bar\rho_K\bar c_K^2,\ \sigma^{\bB}_K=\sqrt{\bar\rho_K}\,\bar c_K,
\]
\cm{$r_{loc}=1$}. Each element is measured against its own scales, so the corona ($\rho\sim10^{-8}\rho_0$) is
not hidden by the photosphere.
\item \textbf{Ratio, cap, coefficient.} \cm{$\mathcal R_K=\max_{I\in K}\max_{s\le8}R^s_I/n^s_K$};
\cm{$w_K=\max_I(|\bv_I|+\sqrt{\max(\gamma p_I/\rho_I+|\bB_I|^2/\rho_I,0)})$};
\[
\nuc=\max\big(0,\min(0.5\,\Delta_Kw_K,\ \Delta_K^2\mathcal R_K)\big),\qquad
\nufl=0.03\,\Delta_Kw_K,\qquad \nu_K=\max(\nuc,\nufl),
\]
\cm{one value per element}, \cm{recomputed at every stage}.
\item \textbf{Slot coefficients} (conserved branch with \code{lnazarov\_energy}, \code{SGS.jl:1881--1885}):
$\nu_K$ on slots $1,2,3,5,\dots,9$; slot 4 gets $\kappa^E_K=\max\big(\gamma(\gamma-1)/\Pr\cdot\nuc,\ \nufl\big)=\max(0.0525\,\nuc,\nufl)$.
\end{enumerate}

\section{The viscous operator}
\code{rhs.jl:1533} sets \code{dsgs\_split\_energy[] = true}; \code{user\_primitives.jl} then returns the
departure from $q_e$ with the energy split into its non-thermal and thermal parts,
$\delta q=q-q_e$, $E_{th}=p/(\gamma-1)$, $E_{nth}=E-E_{th}$, the latter in the spare slot
$n_{eqs}+2=11$ of \code{uprimitive}. \code{\_expansion\_visc!} (\code{rhs.jl:3050--3066}) assembles
\begin{center}\small
\begin{tabular}{@{}lll@{}}
\toprule
Equation & Viscous flux $\bF_V$ & Remark\\
\midrule
$\rho$ & $\nu\nabla\delta\rho$ & \cm{mass diffusion on}\\
$\rho u,\rho v,\rho w$ & $\nu\nabla\delta(\rho u)$, \dots & scalar Laplacians, no stress\\
$E$ & $\nu\nabla\delta E_{nth}+\kappa^E\nabla\delta E_{th}$ & $\nu$ = the $\rho u$-slot coefficient\\
$B_k$ & $\nu\nabla\delta B_k$ & $\eta=\nu$\\
$\psi$ & $\nu\nabla\psi$ & $\psi_e=0$\\
\bottomrule
\end{tabular}
\end{center}
\mj{A Laplacian of the conserved state (here: of its departure from $q_e$)}, in divergence form, with
\mj{no $\tau\cdot\bv$ or resistive-work term}. At rest ($q=q_e$) the operator vanishes.

\section{Positivity}
\subsection*{Case-local floor as the RK stage limiter (\code{fe\_positivity\_limiter!}, \code{user\_inputs.jl:29--77})}
This is \emph{not} the \code{:lpositivity} machinery of \code{src/kernel/positivity} (\code{:lpositivity}
stays \code{false}); it is a function of the deck handed to the solver as
\code{CarpenterKennedy2N54(stage\_limiter! = fe\_positivity\_limiter!, williamson\_condition = false)}
(the keyword since \code{e448f7d}): the floor runs before the RHS of stages 2--5 and after stage 5. (At \commit{} the
keyword was not set and the same behaviour came only from \code{calck = true}, which Jexpresso's \code{solve} implies
by passing a non-empty \code{saveat}.)
Per node $I$, with $f=0.2$ (\code{fe\_floor\_frac}) and the reference values at that node:
\begin{enumerate}[leftmargin=*]
\item if $\rho_I<f\rho_{e,I}$: $\rho_I\leftarrow f\rho_{e,I}$ (momentum unchanged, i.e.\ the velocity is reduced);
\item with $k=|\bmom_I|^2/(2\rho_I)$ (new $\rho$) and $m=\tfrac12|\bB_I|^2+\tfrac12\psi_I^2$: if
$(\gamma-1)(E_I-k-m)<f\,p_{e,I}$, then $E_I\leftarrow f\,p_{e,I}/(\gamma-1)+k+m$ ($\bmom$, $\bB$, $\psi$ unchanged).
\end{enumerate}
It \textbf{adds mass and energy} wherever it fires (not conservative). Hits are counted per rank and
reported at the first hit and each further decade, with a sample $(x,z)$; the count is rank-local
(printed with the \code{MPI.COMM\_WORLD} rank).

\subsection*{Other ingredients}
\begin{itemize}[leftmargin=*]
\item Flux-only floor: $p\leftarrow\max(p,10^{-9})$ inside \code{user\_flux!} (\code{p\_floor\_mhd}); the state is untouched.
\item The \dsgs{} mass diffusion and the $C_{\min}$ background act in the same direction.
\end{itemize}

\section{Observations: code versus comments}
\begin{enumerate}[leftmargin=*]
\item \textbf{[Fixed in \code{f03762b}.]} At \commit{} the deck said ``Slot 4 carries the TOTAL ENERGY: \code{"energy"}
keeps the kernel's $\tau\cdot u$ \dots active'' (under \code{:dsgs\_conserved => true} the term is switched off,
\code{rhs.jl:2965}), ``$\Delta t=2.5\times10^{-3}\tau_0$ is CFL $\approx0.07$'' (the deck sets $7.5\times10^{-3}$:
CFL $\approx0.22$ initially, $\approx0.43$ late), and that the energy slot is ``scaled by $\gamma(\gamma-1)/\Pr$''
(the first implementation; the kernel splits the energy flux). The comments and the README's sizing table now
say what the code does.
\item \textbf{[Fixed in \code{e448f7d}.]} At \commit{} the floor reached stages 2--5 only because \code{calck = true} in
Jexpresso's \code{solve}; the deck did not set \code{williamson\_condition = false} as the conservative limiter does
(\code{positivity\_driver.jl:266--268}), and a solve without callbacks or \code{saveat} would have applied it once per
step (measured on the installed OrdinaryDiffEqLowStorageRK 1.11.0: 1 call per step against 5). The deck now sets the
keyword; through Jexpresso's \code{solve} (\code{calck = true}) this changes nothing: run here from the same
initial condition, the decks of \code{c1d5b46} and \code{e448f7d} give byte-identical VTK output at $t=1$ and
$2\,\tau_0$ ($\Delta t=7.5\times10^{-3}$). The fix only removes the dependence on \code{calck}.
\end{enumerate}

% Shared appendix: where each mechanism lives in the source (branch sm/MHD, commit dcc3653).
\section*{Appendix: code map (branch \branch, commit \commit)}
\addcontentsline{toc}{section}{Appendix: code map}
{\small
\begin{longtable}{@{}L{0.47\textwidth}L{0.47\textwidth}@{}}
\toprule
\textbf{Mechanism} & \textbf{Where (file:line)}\\
\midrule\endhead
\bottomrule\endfoot
Deck defaults of every \code{:dsgs\_*} key & \code{src/io/mod\_inputs.jl:1085--1376}\\
Deck defaults of \code{:lpositivity}, floors, report & \code{src/io/mod\_inputs.jl:456--477}\\
\code{:positivity\_method}; rebuild of the RK solver & \code{src/io/mod\_inputs.jl:828--838}\\
\dsgs{} history rotation, stage weights, hold, freeze & \code{src/kernel/operators/rhs.jl:927--984}\\
Stencil selection \code{\_dsgs\_stencil} & \code{rhs.jl:141}\\
Hold length \code{\_dsgs\_hold\_steps} ($\max(2,H)+1$ rotations) & \code{rhs.jl:182}\\
Residual right-hand side (element / legacy / Dirichlet) & \code{rhs.jl:188--280}; legacy fill \code{rhs.jl:282}; boundary zeroing \code{rhs.jl:313, 346}\\
1D \code{DSGS\_MHD} dispatch & \code{rhs.jl:1398--1470}\\
2D \code{DSGS\_MHD} dispatch (sets \code{dsgs\_ref\_weight}, \code{dsgs\_split\_energy}, \code{dsgs\_nodal\_rho}) & \code{rhs.jl:1517--1610}\\
2D viscous assembly \code{\_viscous\_rhs\_el\_2d\_dsgs!} & \code{rhs.jl:1796}\\
2D weak viscous flux \code{\_expansion\_visc!} (conserved flag, ref.\ weight, split energy, $\tau\cdot\bv$, resistive work) & \code{rhs.jl:2943--3210} (\code{lcons} 2959, weight 2978, split 3050, $\tau\cdot\bv$ 3076, $\eta\,\bB\cdot\nabla\bB$ 3091)\\
1D weak viscous flux & \code{rhs.jl:2823}; element loop \code{rhs.jl:1321}\\
2D \code{DSGS\_MHD} element kernel & \code{src/kernel/physics/SGS.jl:1598--1912}\\
1D \code{DSGS\_MHD} element kernel & \code{SGS.jl:1929--2120}\\
Nodal kernels (off in every MHD deck; BW can switch it on with \code{JEXPRESSO\_BW\_NODAL=1}) & \code{SGS.jl:2144} (1D), \code{SGS.jl:2515} (2D)\\
Normalization floor \code{\_dsgs\_denom}, \code{:dsgs\_rel} multiplier & \code{SGS.jl:676--679}; set per call in \code{rhs.jl:192}\\
Smoothness cutoff \code{\_dsgs\_cut} & \code{SGS.jl:703}\\
\code{SGS\_diffusion(::DSGS\_MHD)} (nodal-$\rho$ switch) & \code{SGS.jl:547--569}\\
Element length $\Delta^{elem}_K$ (2D: min of corner extents) & \code{src/kernel/mesh/mesh.jl:6430--6457}; 1D $\Delta x$: \code{mesh.jl:5863}\\
\code{uprimitive} width ($n_{eqs}+2$ in 2D, $n_{eqs}+1$ in 1D) & \code{src/kernel/globalStructs.jl:139--153}\\
Realizability repair (\code{"repair"}, unused by the MHD decks) & \code{src/kernel/positivity/Positivity.jl:294}; driver \code{positivity\_driver.jl:124}; called in \code{rhs.jl:821}\\
Conservative element limiter (\code{"conservative"}) & \code{Positivity.jl:433--559}; driver \code{positivity\_driver.jl:245--337}\\
Main \code{solve} call (\code{saveat} non-empty $\Rightarrow$ \code{calck = true}) & \code{src/kernel/solvers/TimeIntegrators.jl:855}\\
\end{longtable}}
\noindent\small Stage-limiter calls in \code{CarpenterKennedy2N54}: with \code{williamson\_condition = false}
the limiter is called before the right-hand side of stages 2--5 and once after stage 5
(OrdinaryDiffEq v6.90/v6.95, \code{low\_storage\_rk\_perform\_step.jl}). In OrdinaryDiffEq
\code{calck = true} forces \code{williamson\_condition = false} in the 2N cache, and Jexpresso's
\code{solve} passes a non-empty \code{saveat}, so this holds for every deck that hands the solver
a stage limiter. (Jexpresso ships no \code{Manifest.toml}; the package version is whatever the
environment resolves, and the check above was made against v6.90/v6.95.)

\end{document}
